\chapter{Introduction to Data Mining and the KDD Process}
\label{ch:intro-data-mining}

\section{Definition and Fundamental Motivations}

The constant technological progress in data generation and storage has led to an exponential growth of available information, both in commercial and scientific contexts. The modern paradigm of data collection is often summarized in the principle of gathering any possible data, anywhere and anytime, in the expectation that such data may acquire value not only for the original purpose for which they were collected, but also for secondary purposes not initially foreseen.

\begin{definition}[Data Mining]
\textbf{Data Mining} is the application of efficient techniques and algorithms to the analysis of large data collections, aimed at the extraction of valid, novel, potentially useful, and interpretable patterns (hidden knowledge).
\end{definition}

Data Mining does not consist in the simple execution of queries on traditional databases (for example, retrieving the average balance via SQL), but rather in the automatic or semi-automatic discovery of non-trivial structures, dependencies, or regularities in the data.

\subsection{Application Perspectives}
The motivations behind the adoption of Data Mining techniques can be distinguished primarily according to two points of view:

\begin{itemize}
    \item \textbf{Commercial point of view:} Companies manage enormous volumes of data (\textit{data warehousing}): web data (Yahoo, Facebook), e-commerce transactions and retail sales (Amazon), credit card records and bank accounts. In the face of increasingly cheap and accessible computing power, competitive pressure requires extracting value from such masses of data to improve and personalize the services offered (for example in \textit{Customer Relationship Management} - CRM).
    \item \textbf{Scientific and social point of view:} In disciplines such as astronomy (\textit{Sky Survey Data}), Earth observation (NASA EOSDIS satellites collect petabytes of environmental data per year), genomics (gene expression, proteomics), neuroimaging (fMRI), and complex numerical simulations, data are generated at such speeds as to make manual analysis impossible. Data Mining supports scientists in automatic analysis and in the guided formulation of new hypotheses.
\end{itemize}

Big Data represent to all intents and purposes a proxy of social life (purchasing behaviors, lifestyles, opinions, relationship networks, territorial mobility), offering great opportunities to address global issues: from the efficiency of healthcare to the search for clean energy sources, from forecasting climate impacts to optimizing agricultural production.

\subsection{A Confluence of Disciplines}
Data Mining is not an isolated discipline, but rather a synergistic confluence of several areas of computer science and applied mathematics:

\begin{center}
\begin{tikzpicture}[
    node distance=2.2cm,
    every node/.style={font=\small},
    center_node/.style={circle, draw=blue!70, fill=blue!10, thick, inner sep=6pt, align=center, font=\bfseries\small},
    satellite/.style={rectangle, rounded corners, draw=gray!70, fill=gray!5, thick, inner sep=5pt, align=center}
]
    \node[center_node] (dm) {Data\\Mining};
    \node[satellite, above left=1.2cm and 1.8cm of dm] (stat) {Statistics};
    \node[satellite, above right=1.2cm and 1.8cm of dm] (db) {Database\\Systems};
    \node[satellite, below left=1.2cm and 1.8cm of dm] (ml) {Machine\\Learning};
    \node[satellite, below right=1.2cm and 1.8cm of dm] (alg) {Algorithm\\Theory};
    \node[satellite, above=1.5cm of dm] (vis) {Visualization};
    \node[satellite, below=1.5cm of dm] (other) {Other Disciplines\\(HPC, Optimization)};

    \draw[thick, -Stealth, color=blue!60] (stat) -- (dm);
    \draw[thick, -Stealth, color=blue!60] (db) -- (dm);
    \draw[thick, -Stealth, color=blue!60] (ml) -- (dm);
    \draw[thick, -Stealth, color=blue!60] (alg) -- (dm);
    \draw[thick, -Stealth, color=blue!60] (vis) -- (dm);
    \draw[thick, -Stealth, color=blue!60] (other) -- (dm);
\end{tikzpicture}
\end{center}

\section{Types of Data Origin: Primary and Secondary}
In setting up an analytical project, it is essential to distinguish the origin of the analyzed data:
\begin{itemize}
    \item \textbf{Primary Data:} Original data collected specifically for a precise investigation objective. They have not undergone any preventive alterations or manipulations by third parties.
    \item \textbf{Secondary Data:} Data that had already been previously collected for operational purposes or different from the current one and made available to the analyst. They often derive from the merging of heterogeneous sources and require deep normalization and quality verification operations.
\end{itemize}

\section{The KDD (Knowledge Discovery in Databases) Process}
Data Mining constitutes a specific, although central, phase within the broader knowledge discovery process, known as \textbf{KDD} (\textit{Knowledge Discovery in Databases}). The KDD is an iterative and interactive process composed of several sequential phases with feedback loops:

\begin{center}
\begin{tikzpicture}[
    node distance=1.1cm,
    font=\small,
    box/.style={rectangle, rounded corners, draw=blue!80, fill=blue!5, thick, minimum width=4.5cm, minimum height=0.8cm, align=center},
    data/.style={cylinder, shape border rotate=90, draw=gray!80, fill=gray!10, aspect=0.25, minimum height=0.9cm, minimum width=3.2cm, align=center}
]
    \node[data] (db) {Heterogeneous Sources\\(Databases, Logs, Web)};
    \node[box, below=0.8cm of db] (clean) {1. Data Cleaning \& Integration};
    \node[data, below=0.8cm of clean] (dw) {Data Warehouse};
    \node[box, below=0.8cm of dw] (select) {2. Data Selection \& Transformation};
    \node[data, below=0.8cm of select] (taskdata) {Task-relevant Data};
    \node[box, below=0.8cm of taskdata, draw=red!80, fill=red!5] (dm) {3. Data Mining (Algorithms)};
    \node[box, below=0.8cm of dm] (eval) {4. Pattern Evaluation \&\\Knowledge Presentation};

    \draw[-Stealth, thick] (db) -- (clean);
    \draw[-Stealth, thick] (clean) -- (dw);
    \draw[-Stealth, thick] (dw) -- (select);
    \draw[-Stealth, thick] (select) -- (taskdata);
    \draw[-Stealth, thick] (taskdata) -- (dm);
    \draw[-Stealth, thick] (dm) -- (eval);

    \draw[-Stealth, dashed, thick, color=gray!80] (eval.east) to [bend right=45] node[right, font=\scriptsize] {Iterative refinement} (select.east);
\end{tikzpicture}
\end{center}

Each phase performs specific tasks:
\begin{enumerate}
    \item \textbf{Business Understanding:} Understanding the domain objectives and translating the business problem into a formal data analysis problem.
    \item \textbf{Data Cleaning and Data Integration:} 
    \begin{itemize}
        \item \textit{Data Cleaning}: removal of noise, correction of syntactic errors, and handling of inconsistencies.
        \item \textit{Data Integration}: aggregation and reconciliation of data from multiple and heterogeneous sources, typically flowing into a \textbf{Data Warehouse} (a structured and optimized repository for analysis and decision support).
    \end{itemize}
    \item \textbf{Data Selection and Data Transformation:}
    \begin{itemize}
        \item \textit{Data Selection}: extraction of the subset of attributes and records strictly relevant to the investigated task (\textit{task-relevant data}).
        \item \textit{Data Transformation}: transformation of features through aggregations, rescalings, normalizations, or discretizations so that they are suitable for the algorithm.
    \end{itemize}
    \item \textbf{Data Mining:} The actual execution of intelligent algorithms for the extraction of patterns and mathematical models.
    \item \textbf{Pattern Evaluation and Knowledge Presentation:} Identification of truly significant patterns through objective and subjective metrics, elimination of redundancy, and graphical visualization of results using tools interpretable by the end user.
\end{enumerate}

\section{Taxonomy of Data Mining Tasks}

Following the classification of Fayyad et al. \cite{fayyad1996}, Data Mining tasks are divided into two macro-categories:
\begin{itemize}
    \item \textbf{Prediction Methods:} Aim to estimate the future or unknown value of a specific target variable starting from the values of other observed variables (predictor attributes).
    \item \textbf{Description Methods:} Aim to identify human-readable and interpretable patterns that describe the intrinsic structure and hidden relationships in the data.
\end{itemize}

\subsection{Predictive Tasks: Classification}
\textbf{Classification} aims to learn an objective function $f$ that maps each attribute record $x$ into a specific discrete class label $y \in \mathcal{Y}$.

\begin{figure}[htbp]
\centering
\begin{tikzpicture}[font=\small, node distance=1.5cm]
    \node[rectangle, draw=blue!70, fill=blue!5, thick, text width=3.2cm, align=center] (train) {\textbf{Training Set}\\Historical records with known class $y$};
    \node[rectangle, draw=red!70, fill=red!5, thick, text width=2.8cm, align=center, right=1.2cm of train] (learn) {\textbf{Learning Algorithm}\\(e.g., Decision Tree)};
    \node[rectangle, draw=green!70, fill=green!5, thick, text width=2.6cm, align=center, right=1.2cm of learn] (model) {\textbf{Classifier Model}\\$y = f(x)$};
    \node[rectangle, draw=gray!70, fill=gray!5, thick, text width=3.2cm, align=center, below=1.2cm of model] (test) {\textbf{Test Set / New Data}\\Records with class to be estimated};
    \node[rectangle, draw=purple!70, fill=purple!5, thick, text width=2.6cm, align=center, left=1.2cm of test] (pred) {\textbf{Prediction}\\Class label};

    \draw[-Stealth, thick] (train) -- (learn);
    \draw[-Stealth, thick] (learn) -- (model);
    \draw[-Stealth, thick] (model) -- (test);
    \draw[-Stealth, thick] (test) -- (pred);
\end{tikzpicture}
\caption{Conceptual workflow of the classification process.}
\label{fig:classification-workflow}
\end{figure}

The dataset is divided into two disjoint sets:
\begin{itemize}
    \item \textbf{Training Set:} Used to train the classifier.
    \item \textbf{Test Set:} Used to evaluate the generalization capability of the model on unseen instances during the learning phase.
\end{itemize}

\begin{example}[Classification Applications]
Among the main and most widespread application areas of classification models we find:
\begin{itemize}
    \item \textbf{Credit Worthiness Evaluation:} Predicting whether a loan applicant will repay the debt (\textit{Yes/No}) based on employment, education level, and years of residence.
    \item \textbf{Bank Fraud Detection:} Classifying credit card transactions as legitimate or fraudulent. Features include amount, time, geographical location, and historical spending profile.
    \item \textbf{Churn Prediction:} Estimating the probability that a telecommunications customer will abandon the service for a competitor, analyzing call duration, data traffic volume, and recorded complaints.
    \item \textbf{Other applications:} Automatic categorization of news (sports, finance, politics), clinical diagnosis (benign or malignant tumor), land cover classification from satellite images.
\end{itemize}
\end{example}

\subsection{Descriptive Tasks}

\subsubsection{1. Clustering}
\textbf{Clustering} consists in organizing a set of objects into distinct groups (called \textit{clusters}), such that:
\begin{enumerate}
    \item The distance between objects belonging to the same cluster (\textbf{intra-cluster distance}) is \textbf{minimized} (maximum internal homogeneity).
    \item The distance between objects belonging to different clusters (\textbf{inter-cluster distance}) is \textbf{maximized} (maximum external separation).
\end{enumerate}

\begin{figure}[htbp]
\centering
\begin{tikzpicture}[scale=0.85]
    % Cluster 1
    \draw[thick, fill=blue!10, draw=blue!60, rounded corners=15pt] (-3.5,-1) rectangle (-1.2,1.5);
    \fill[blue!70] (-3,0.8) circle (3pt);
    \fill[blue!70] (-2.2,1.1) circle (3pt);
    \fill[blue!70] (-2.7,0.1) circle (3pt);
    \fill[blue!70] (-1.8,-0.4) circle (3pt);
    \node[blue!80!black, font=\bfseries\scriptsize] at (-2.35,-0.7) {Cluster 1};

    % Cluster 2
    \draw[thick, fill=green!10, draw=green!60, rounded corners=15pt] (1.2,-1) rectangle (3.5,1.5);
    \fill[green!60!black] (1.7,0.9) circle (3pt);
    \fill[green!60!black] (2.4,1.1) circle (3pt);
    \fill[green!60!black] (3.0,0.2) circle (3pt);
    \fill[green!60!black] (2.1,-0.5) circle (3pt);
    \node[green!60!black, font=\bfseries\scriptsize] at (2.35,-0.7) {Cluster 2};

    % Intra-cluster distance arrow
    \draw[Stealth-Stealth, thick, blue!80] (-2.7,0.1) -- node[above, font=\tiny] {min intra-distance} (-1.8,-0.4);

    % Inter-cluster distance arrow
    \draw[Stealth-Stealth, very thick, red!80] (-1.2,0.5) -- node[above, font=\scriptsize\bfseries] {max inter-distance} (1.2,0.5);
\end{tikzpicture}
\caption{Fundamental geometric principle of clustering.}
\label{fig:clustering-principle}
\end{figure}

Clustering is an unsupervised technique: it does not require a priori classes.
\begin{itemize}
    \item \textbf{Market Segmentation:} Grouping customers based on socio-demographic data and purchasing habits to target differentiated marketing campaigns.
    \item \textbf{Document Clustering:} Grouping similar textual documents based on the joint frequencies of keywords (for example, the analysis of the Enron corporate email archive).
    \item \textbf{Summarization and Earth Sciences:} Reducing the complexity of massive spatial series by partitioning sea surface temperature (\textit{Sea Surface Temperature} - SST) and primary production (\textit{Net Primary Production} - NPP) measurements in the Northern and Southern hemispheres.
\end{itemize}

\subsubsection{2. Association Analysis (Discovery of Association Rules)}
Given a set of transactions, each consisting of a subset of items taken from a global set of items, association analysis aims to discover probabilistic dependencies in the form of rules:
\begin{equation}
X \implies Y
\end{equation}
where $X$ and $Y$ are disjoint itemsets ($X \cap Y = \emptyset$). The rule indicates that the presence of the items in the antecedent $X$ systematically predicts the contextual presence of the item or items in the consequent $Y$.

\begin{example}[Market Basket Analysis]
Considering a history of cash register receipts:
\begin{table}[htbp]
\centering
\begin{tabular}{cl}
\toprule
\textbf{TID} & \textbf{Purchased Items} \\
\midrule
1 & Bread, Coca-Cola, Milk \\
2 & Beer, Bread \\
3 & Beer, Coca-Cola, Diapers, Milk \\
4 & Beer, Bread, Diapers, Milk \\
5 & Coca-Cola, Diapers, Milk \\
\bottomrule
\end{tabular}
\end{table}
The algorithm extracts rules with high statistical evidence, such as:
\[
\{\text{Milk}\} \implies \{\text{Coca-Cola}\}, \quad \{\text{Diapers}, \text{Milk}\} \implies \{\text{Beer}\}
\]
Such information guides the strategic placement of goods on shelves and the design of cross-selling offers.
\end{example}

\textbf{Other applications:} Diagnosis of alarms in telecommunication networks (identifying temporal patterns of network alarms that trigger jointly to identify systemic root failures); medical informatics (correlation between combinations of clinical symptoms, blood tests, and emerging pathologies).

\subsubsection{3. Anomaly / Deviation / Change Detection}
It consists in identifying objects that deviate significantly from the average or expected behavior of the population. 
Anomalies can be generated by instrumental failures or entry errors (and in this case represent \textit{noise} to be purged), or they can reflect exceptional events of extreme application interest (credit fraud, cyber intrusions against corporate networks, monitoring of sudden deforestation of forest cover via remote sensors).

\section{Methodological and Computational Challenges in Data Mining}
The application of traditional statistical analysis and database methods is often inadequate due to five primary challenges:
\begin{enumerate}
    \item \textbf{Scalability:} Algorithms must handle data volumes on the order of terabytes or petabytes without computational complexity and memory footprint exploding.
    \item \textbf{High Dimensionality:} The presence of hundreds or thousands of attributes exposes one to the so-called \textit{curse of dimensionality}, where the geometric space becomes extremely empty and distances tend to lose discriminative power.
    \item \textbf{Complex and Heterogeneous Data:} Need to process semi-structured or unstructured structures (graphs, texts, genomic sequences, spatiotemporal series).
    \item \textbf{Data Ownership and Distribution:} Data are frequently geographically distributed across multiple organizations and subject to strict legal privacy and security constraints.
    \item \textbf{Non-Traditional Analysis:} Need to extract knowledge from continuous streams in real-time (\textit{data stream mining}) without the possibility of storing the entire dataset.
\end{enumerate}
